% ============================================================
%  插图绘制参考（模板文件，默认不编入设计说明书 PDF）
%  设计说明书正文只保留 figures/ 中实际出图的 18 张插图：
%    界面截图 10 张、PlantUML 结构图 8 张（图号由正文自动连续编号）。
%  本文件保留这 8 张结构图的可编译 PlantUML 源码与 10 张截图的取图步骤，
%  供后续补图、改图时参考；如需把本章编入 PDF，在 design.tex 的
%  \end{document} 前加入 \appendix 与 % ============================================================
%  插图绘制参考（模板文件，默认不编入设计说明书 PDF）
%  设计说明书正文只保留 figures/ 中实际出图的 18 张插图：
%    界面截图 10 张、PlantUML 结构图 8 张（图号由正文自动连续编号）。
%  本文件保留这 8 张结构图的可编译 PlantUML 源码与 10 张截图的取图步骤，
%  供后续补图、改图时参考；如需把本章编入 PDF，在 design.tex 的
%  \end{document} 前加入 \appendix 与 % ============================================================
%  插图绘制参考（模板文件，默认不编入设计说明书 PDF）
%  设计说明书正文只保留 figures/ 中实际出图的 18 张插图：
%    界面截图 10 张、PlantUML 结构图 8 张（图号由正文自动连续编号）。
%  本文件保留这 8 张结构图的可编译 PlantUML 源码与 10 张截图的取图步骤，
%  供后续补图、改图时参考；如需把本章编入 PDF，在 design.tex 的
%  \end{document} 前加入 \appendix 与 % ============================================================
%  插图绘制参考（模板文件，默认不编入设计说明书 PDF）
%  设计说明书正文只保留 figures/ 中实际出图的 18 张插图：
%    界面截图 10 张、PlantUML 结构图 8 张（图号由正文自动连续编号）。
%  本文件保留这 8 张结构图的可编译 PlantUML 源码与 10 张截图的取图步骤，
%  供后续补图、改图时参考；如需把本章编入 PDF，在 design.tex 的
%  \end{document} 前加入 \appendix 与 \input{chapters/figure_spec}。
% ============================================================

\chapter{插图绘制参考}

本说明书的插图统一放在 \code{figures/} 目录，文件名为小写英文与下划线；
结构图由 PlantUML 渲染为 PNG，界面截图在运行的系统中按路由截取。
图号、图题与交叉引用均由正文 \verb|\figph| 生成，增删图片后自动重排。

\section{通用出图约定}

\begin{itemize}[leftmargin=2em,itemsep=0.2em]
\item 结构图：复制下面的 \code{@startuml ... @enduml} 片段保存为同名 \code{.puml}，
  执行 \code{plantuml -tpng figures/<文件名去掉.png>.puml} 生成 PNG；中文渲染需本机中文字体。
\item 界面截图：启动 \code{python run.py --no-browser}，用 Chrome/Edge 在 1440×900 及以上分辨率截取 PNG。
\item 图片宽度建议不小于 1400 像素；正文按版心宽度 92\% 缩放，最大高度为版心高度 70\%。
\end{itemize}

\section{图~\ref{fig:arch01}：系统总体架构图}

文件 \code{arch_01_system_overview.png}，组件图。画法：按七层自上而下排列，前端三块（页面组、回放、仿真）→ 接口层三个组件 → 业务服务层五个组件 → 独立的仿真子系统包 → 适配器层 → 数据存储层 → 外部后端；用实线箭头标注 HTTP/JSON、SSE、MJPEG、子进程、文件读写等协议。

\begin{lstlisting}[style=navplain]
@startuml arch_01_system_overview
skinparam componentStyle rectangle
skinparam defaultFontName "Noto Sans CJK SC"
title 码极客 NavEval Studio V1.0 系统总体架构图
actor "用户 admin/operator/viewer" as User

package "浏览器前端 Vue 3" {
  [总览/数据/任务/分析页面组] as UI1
  [回放工作台 Canvas] as UI2
  [交互式仿真 MJPEG] as UI3
}
package "接口服务层 FastAPI" {
  [REST 路由 / Pydantic 校验] as API
  [Bearer + ?token 鉴权] as AUTH
  [SSE 事件流] as SSE
}
package "业务服务层" {
  [数据扫描/切分/质量校验] as SVC1
  [任务编排/调度器] as SVC2
  [指标聚合/归因/报告] as SVC3
  [回放解析/资产导入] as SVC4
  [后台守护 自动同步] as SVC5
}
package "仿真子系统（可选）" {
  [会话管理/动作队列] as SIM1
  [引擎线程 渲染循环] as SIM2
  [策略/录制/俯视绘制] as SIM3
}
package "适配器层" {
  [BaseAdapter 协议] as ADP
  [StreamVLNAdapter] as ADP1
  [MockAdapter] as ADP2
}
package "数据存储层" {
  [SQLite WAL eval.db] as DB
  [工作目录 storage/] as FS
}
package "外部推理与仿真后端" {
  [StreamVLN 推理子进程] as EXT1
  [habitat_sim 场景渲染] as EXT2
}
User --> UI1
UI1 --> API : HTTP / JSON
UI2 --> API : 回放 JSON + 帧图像
UI3 --> SSE : 事件流
UI3 --> API : MJPEG / ?token=
API --> AUTH
API --> SVC1
API --> SVC2
API --> SVC3
API --> SVC4
SVC2 --> SSE : 进度/采样/日志
SVC2 --> ADP
ADP <|-- ADP1
ADP <|-- ADP2
ADP1 --> EXT1 : 子进程 + --dump_frames
SVC4 --> SIM1
SIM1 --> SIM2
SIM1 --> SIM3
SIM2 --> EXT2
SIM3 --> EXT2
SVC1 --> DB
SVC2 --> DB
SVC3 --> DB
SVC4 --> DB
SVC5 --> DB
SVC5 --> FS : 备份库目录监视
SVC2 --> FS : 快照/日志/逐帧 jsonl
SIM3 --> FS : 录制帧与 mp4
@enduml
\end{lstlisting}
\section{图~\ref{fig:flow01}：评测任务执行总流程}

文件 \code{flow_01_eval_task.png}，活动图，软著材料的核心流程图。画法：用泳道区分前端、接口服务层、调度器、适配器/子进程、数据存储，覆盖提交、校验、排队、执行、事件回写、聚合六个阶段。

\begin{lstlisting}[style=navplain]
@startuml flow_01_eval_task
skinparam defaultFontName "Noto Sans CJK SC"
title 评测任务执行总流程
|前端|
start
:选择数据集/切分/模型/模板并覆盖参数;
:提交创建任务;
|接口服务层|
:参数校验 + 因果性白名单预检;
if (校验通过?) then (否)
  :返回 1001 / 1006 与具体字段;
  |前端|
  :就地提示并保留表单;
  stop
else (是)
endif
|数据存储|
:写入 eval_task(queued) 与配置快照;
|调度器|
:入队，等待并发额度与 GPU 互斥;
:取出任务，创建工作目录;
:写 config.json 与启动命令;
|适配器/子进程|
:prepare 校验路径与后端能力;
:拉起推理子进程（可选 --dump_frames）;
repeat
  :逐条执行 episode;
  :产出进度/资源/日志/逐帧记录;
  |数据存储|
  :按行解析写 task_episode 与 frames/*.jsonl;
  |前端|
  :SSE 更新进度条/曲线/日志;
repeat while (还有 episode?) is (是)
|接口服务层|
:聚合指标、生成回放索引与报告索引;
:任务置为 succeeded，推送 done 事件;
|前端|
:展示结果，支持对比/回放/导出;
stop
@enduml
\end{lstlisting}
\section{图~\ref{fig:flow02}：数据与配置管理子系统流程}

文件 \code{flow_02_data_manage.png}。画法：一条主干（登记→扫描→索引→校验→切分→导出），在校验处分支展示问题清单，在导出处强调时间前缀二次校验。

\begin{lstlisting}[style=navplain]
@startuml flow_02_data_manage
skinparam defaultFontName "Noto Sans CJK SC"
title 数据与配置管理子系统流程
start
:登记数据集并校验根路径白名单;
if (路径合法且存在?) then (否)
  :状态置为 missing 并提示修正;
  stop
else (是)
endif
:扫描 episode 并写入索引，状态置为 ready;
:多源观测对齐与质量校验;
if (存在问题?) then (是)
  :输出问题清单与涉及 episode;
else (否)
  :标记校验通过;
endif
:生成 full / smoke / subset 切分并保存条件;
:因果切片导出与时间前缀二次校验;
:写出 jsonl 并返回下载链接;
stop
@enduml
\end{lstlisting}
\section{图~\ref{fig:flow03}：任务调度子系统流程}

文件 \code{flow_03_scheduler.png}。画法：循环从队列取任务，突出并发额度、GPU 互斥、episode 边界上的暂停/取消检查与失败重跑。

\begin{lstlisting}[style=navplain]
@startuml flow_03_scheduler
skinparam defaultFontName "Noto Sans CJK SC"
title 任务调度子系统流程
start
while (队列非空?) is (是)
  :按优先级取队首任务;
  if (并发额度未满且 GPU 可用?) then (否)
    :留在队列等待下一轮;
  else (是)
    :标记 running 并占用 GPU;
    :prepare 调用适配器;
    while (还有 episode?) is (是)
      :执行一条 episode，写结果并推送事件;
      if (收到暂停?) then (是)
        :在 episode 边界置 paused 并等待恢复;
      else (否)
      endif
      if (收到取消?) then (是)
        :终止子进程，置 cancelled 并释放资源;
        stop
      else (否)
      endif
    endwhile
    :聚合指标并释放并发额度与 GPU;
  endif
endwhile (否)
:队列空闲，等待新任务或重跑;
stop
@enduml
\end{lstlisting}
\section{图~\ref{fig:flow04}：推理适配器与逐帧回放流程}

文件 \code{flow_04_adapter_replay.png}。画法：从 prepare 探测上游能力开始，分“支持/不支持逐帧”两支，最终汇合到回放接口与四区工作台。

\begin{lstlisting}[style=navplain]
@startuml flow_04_adapter_replay
skinparam defaultFontName "Noto Sans CJK SC"
title 推理适配器与逐帧回放流程
start
:读取配置快照，prepare 校验路径与权重;
if (上游支持逐帧落盘?) then (是)
  :追加逐帧参数，上游每步写位姿与记忆状态;
  :解析 jsonl，写逐帧文件并登记 traj_path;
else (否)
  :降级为仅汇总指标;
endif
:汇总 SR / SPL / NE / nDTW 与资源峰值;
:回放接口读取任务号与行标识;
if (存在逐帧文件?) then (是)
  :还原俯视轨迹、第一视角、动作时间轴与记忆曲线;
else (否)
  :显示汇总指标与参考路径;
endif
stop
@enduml
\end{lstlisting}
\section{图~\ref{fig:flow05}：指标分析与诊断流程}

文件 \code{flow_05_analysis.png}。画法：从逐条结果到四条分析支路（对比、消融、显著性、长时曲线）与失败归因，最后汇入报告导出。

\begin{lstlisting}[style=navplain]
@startuml flow_05_analysis
skinparam defaultFontName "Noto Sans CJK SC"
title 指标分析与诊断流程
start
:读取逐条结果并聚合 SR / SPL / NE / nDTW;
fork
  :多任务对比矩阵与基线差值;
fork again
  :模态组合 × 记忆策略消融矩阵;
fork again
  :配对 bootstrap 与 95% 置信区间;
fork again
  :按步数分桶的长时曲线;
end fork
:九类失败归因自动分类与人工复核;
:生成 Markdown / CSV 报告并下载;
stop
@enduml
\end{lstlisting}
\section{图~\ref{fig:flow06}：交互式仿真子系统流程}

文件 \code{flow_06_sim.png}。画法：会话启动后 \code{fork} 出引擎循环、驾驶指令、录制三条并行支路，用 \code{end fork} 汇合到停止与资源释放。

\begin{lstlisting}[style=navplain]
@startuml flow_06_sim
skinparam defaultFontName "Noto Sans CJK SC"
title 交互式仿真子系统流程
start
:读取场景、起点、目标与参考路径;
:单会话检查并启动 Habitat 引擎线程;
fork
  :按节拍渲染第一视角与第三视角画中画;
  :MJPEG 分发并绘制俯视三线轨迹;
fork again
  :手动 / oracle / 模型动作写入队列并按节拍消费;
fork again
  :按 record_fps 采样 JPEG，停止后编码 mp4;
end fork
:输出步数、距离、NE 与 SPL 状态快照;
:空闲超时或用户停止后，停止引擎并释放资源;
stop
@enduml
\end{lstlisting}
\section{图~\ref{fig:flow07}：资产同步导入流程}

文件 \code{flow_07_backup_import.png}。画法：扫描→五类分类→签名比对→dry\_run 分支→写入→运行中结果增量同步。

\begin{lstlisting}[style=navplain]
@startuml flow_07_backup_import
skinparam defaultFontName "Noto Sans CJK SC"
title 资产同步导入流程
start
:读取绑定目录与自动同步开关;
:按五类扫描，排除大目录与超阈值路径;
:计算 source_key 与 signature;
if (与账本一致?) then (是)
  :跳过并保留已有记录;
else (否)
  if (dry_run 预览?) then (是)
    :只输出候选清单，不写库;
  else (否)
    :分类落库：数据集、资产索引、模型档案、历史任务;
  endif
endif
if (结果文件仍在写入?) then (是)
  :保持运行中并按条数增量刷新;
else (否)
  :转为已完成并补齐指标;
endif
:后台守护按间隔监视目录变化;
stop
@enduml
\end{lstlisting}
\section{界面截图：取图步骤}

\scriptsize
\begin{longtable}{@{}p{1.0cm}p{4.3cm}p{2.8cm}p{6.1cm}@{}}
\toprule
图号 & 文件名 & 路由 / 页面 & 取图步骤与画面要素 \\
\midrule
\endfirsthead
\toprule
图号 & 文件名 & 路由 / 页面 & 取图步骤与画面要素 \\
\midrule
\endhead
\bottomrule
\endfoot
\ref{fig:ui01} & \code{ui_01_overview.png} & \code{/overview} & 登录后首页；统计卡、队列、最近任务、趋势图与快捷入口 \\
\ref{fig:ui02} & \code{ui_02_dataset_list.png} & \code{/datasets} & 数据集列表；登记/扫描结果与场景、episode 统计 \\
\ref{fig:ui03} & \code{ui_03_observation_timeline.png} & \code{/datasets/:id} & 某 episode 的 RGB/深度/位姿/动作四轨时间轴与偏差标记 \\
\ref{fig:ui06} & \code{ui_06_task_wizard.png} & \code{/tasks} & 新建评测向导；数据集/切分/模型/模板选择与参数覆盖 \\
\ref{fig:ui07} & \code{ui_07_task_monitor.png} & \code{/tasks/:id} & 运行中任务；进度、当前 episode、日志与显存/延迟曲线 \\
\ref{fig:ui12} & \code{ui_12_analysis_compare.png} & \code{/analysis} & 选 2--3 个任务；指标矩阵与相对基线的差值列 \\
\ref{fig:ui14} & \code{ui_14_bootstrap.png} & \code{/analysis} & 两任务配对 bootstrap；均值差、95\% 置信区间与显著性 \\
\ref{fig:ui18} & \code{ui_18_sim_workbench.png} & \code{/sim} & 启动会话并走若干步；第一视角、俯视轨迹、控制与读数 \\
\ref{fig:ui19} & \code{ui_19_sim_policy.png} & \code{/sim} & 模型驾驶面板；加载/启用/暂停/停止与显存提示 \\
\ref{fig:ui23} & \code{ui_23_backup_import.png} & \code{/backup-import} & 绑定目录总览；五类候选数、扫描预览与自动同步开关 \\
\end{longtable}
\normalsize
。
% ============================================================

\chapter{插图绘制参考}

本说明书的插图统一放在 \code{figures/} 目录，文件名为小写英文与下划线；
结构图由 PlantUML 渲染为 PNG，界面截图在运行的系统中按路由截取。
图号、图题与交叉引用均由正文 \verb|\figph| 生成，增删图片后自动重排。

\section{通用出图约定}

\begin{itemize}[leftmargin=2em,itemsep=0.2em]
\item 结构图：复制下面的 \code{@startuml ... @enduml} 片段保存为同名 \code{.puml}，
  执行 \code{plantuml -tpng figures/<文件名去掉.png>.puml} 生成 PNG；中文渲染需本机中文字体。
\item 界面截图：启动 \code{python run.py --no-browser}，用 Chrome/Edge 在 1440×900 及以上分辨率截取 PNG。
\item 图片宽度建议不小于 1400 像素；正文按版心宽度 92\% 缩放，最大高度为版心高度 70\%。
\end{itemize}

\section{图~\ref{fig:arch01}：系统总体架构图}

文件 \code{arch_01_system_overview.png}，组件图。画法：按七层自上而下排列，前端三块（页面组、回放、仿真）→ 接口层三个组件 → 业务服务层五个组件 → 独立的仿真子系统包 → 适配器层 → 数据存储层 → 外部后端；用实线箭头标注 HTTP/JSON、SSE、MJPEG、子进程、文件读写等协议。

\begin{lstlisting}[style=navplain]
@startuml arch_01_system_overview
skinparam componentStyle rectangle
skinparam defaultFontName "Noto Sans CJK SC"
title 码极客 NavEval Studio V1.0 系统总体架构图
actor "用户 admin/operator/viewer" as User

package "浏览器前端 Vue 3" {
  [总览/数据/任务/分析页面组] as UI1
  [回放工作台 Canvas] as UI2
  [交互式仿真 MJPEG] as UI3
}
package "接口服务层 FastAPI" {
  [REST 路由 / Pydantic 校验] as API
  [Bearer + ?token 鉴权] as AUTH
  [SSE 事件流] as SSE
}
package "业务服务层" {
  [数据扫描/切分/质量校验] as SVC1
  [任务编排/调度器] as SVC2
  [指标聚合/归因/报告] as SVC3
  [回放解析/资产导入] as SVC4
  [后台守护 自动同步] as SVC5
}
package "仿真子系统（可选）" {
  [会话管理/动作队列] as SIM1
  [引擎线程 渲染循环] as SIM2
  [策略/录制/俯视绘制] as SIM3
}
package "适配器层" {
  [BaseAdapter 协议] as ADP
  [StreamVLNAdapter] as ADP1
  [MockAdapter] as ADP2
}
package "数据存储层" {
  [SQLite WAL eval.db] as DB
  [工作目录 storage/] as FS
}
package "外部推理与仿真后端" {
  [StreamVLN 推理子进程] as EXT1
  [habitat_sim 场景渲染] as EXT2
}
User --> UI1
UI1 --> API : HTTP / JSON
UI2 --> API : 回放 JSON + 帧图像
UI3 --> SSE : 事件流
UI3 --> API : MJPEG / ?token=
API --> AUTH
API --> SVC1
API --> SVC2
API --> SVC3
API --> SVC4
SVC2 --> SSE : 进度/采样/日志
SVC2 --> ADP
ADP <|-- ADP1
ADP <|-- ADP2
ADP1 --> EXT1 : 子进程 + --dump_frames
SVC4 --> SIM1
SIM1 --> SIM2
SIM1 --> SIM3
SIM2 --> EXT2
SIM3 --> EXT2
SVC1 --> DB
SVC2 --> DB
SVC3 --> DB
SVC4 --> DB
SVC5 --> DB
SVC5 --> FS : 备份库目录监视
SVC2 --> FS : 快照/日志/逐帧 jsonl
SIM3 --> FS : 录制帧与 mp4
@enduml
\end{lstlisting}
\section{图~\ref{fig:flow01}：评测任务执行总流程}

文件 \code{flow_01_eval_task.png}，活动图，软著材料的核心流程图。画法：用泳道区分前端、接口服务层、调度器、适配器/子进程、数据存储，覆盖提交、校验、排队、执行、事件回写、聚合六个阶段。

\begin{lstlisting}[style=navplain]
@startuml flow_01_eval_task
skinparam defaultFontName "Noto Sans CJK SC"
title 评测任务执行总流程
|前端|
start
:选择数据集/切分/模型/模板并覆盖参数;
:提交创建任务;
|接口服务层|
:参数校验 + 因果性白名单预检;
if (校验通过?) then (否)
  :返回 1001 / 1006 与具体字段;
  |前端|
  :就地提示并保留表单;
  stop
else (是)
endif
|数据存储|
:写入 eval_task(queued) 与配置快照;
|调度器|
:入队，等待并发额度与 GPU 互斥;
:取出任务，创建工作目录;
:写 config.json 与启动命令;
|适配器/子进程|
:prepare 校验路径与后端能力;
:拉起推理子进程（可选 --dump_frames）;
repeat
  :逐条执行 episode;
  :产出进度/资源/日志/逐帧记录;
  |数据存储|
  :按行解析写 task_episode 与 frames/*.jsonl;
  |前端|
  :SSE 更新进度条/曲线/日志;
repeat while (还有 episode?) is (是)
|接口服务层|
:聚合指标、生成回放索引与报告索引;
:任务置为 succeeded，推送 done 事件;
|前端|
:展示结果，支持对比/回放/导出;
stop
@enduml
\end{lstlisting}
\section{图~\ref{fig:flow02}：数据与配置管理子系统流程}

文件 \code{flow_02_data_manage.png}。画法：一条主干（登记→扫描→索引→校验→切分→导出），在校验处分支展示问题清单，在导出处强调时间前缀二次校验。

\begin{lstlisting}[style=navplain]
@startuml flow_02_data_manage
skinparam defaultFontName "Noto Sans CJK SC"
title 数据与配置管理子系统流程
start
:登记数据集并校验根路径白名单;
if (路径合法且存在?) then (否)
  :状态置为 missing 并提示修正;
  stop
else (是)
endif
:扫描 episode 并写入索引，状态置为 ready;
:多源观测对齐与质量校验;
if (存在问题?) then (是)
  :输出问题清单与涉及 episode;
else (否)
  :标记校验通过;
endif
:生成 full / smoke / subset 切分并保存条件;
:因果切片导出与时间前缀二次校验;
:写出 jsonl 并返回下载链接;
stop
@enduml
\end{lstlisting}
\section{图~\ref{fig:flow03}：任务调度子系统流程}

文件 \code{flow_03_scheduler.png}。画法：循环从队列取任务，突出并发额度、GPU 互斥、episode 边界上的暂停/取消检查与失败重跑。

\begin{lstlisting}[style=navplain]
@startuml flow_03_scheduler
skinparam defaultFontName "Noto Sans CJK SC"
title 任务调度子系统流程
start
while (队列非空?) is (是)
  :按优先级取队首任务;
  if (并发额度未满且 GPU 可用?) then (否)
    :留在队列等待下一轮;
  else (是)
    :标记 running 并占用 GPU;
    :prepare 调用适配器;
    while (还有 episode?) is (是)
      :执行一条 episode，写结果并推送事件;
      if (收到暂停?) then (是)
        :在 episode 边界置 paused 并等待恢复;
      else (否)
      endif
      if (收到取消?) then (是)
        :终止子进程，置 cancelled 并释放资源;
        stop
      else (否)
      endif
    endwhile
    :聚合指标并释放并发额度与 GPU;
  endif
endwhile (否)
:队列空闲，等待新任务或重跑;
stop
@enduml
\end{lstlisting}
\section{图~\ref{fig:flow04}：推理适配器与逐帧回放流程}

文件 \code{flow_04_adapter_replay.png}。画法：从 prepare 探测上游能力开始，分“支持/不支持逐帧”两支，最终汇合到回放接口与四区工作台。

\begin{lstlisting}[style=navplain]
@startuml flow_04_adapter_replay
skinparam defaultFontName "Noto Sans CJK SC"
title 推理适配器与逐帧回放流程
start
:读取配置快照，prepare 校验路径与权重;
if (上游支持逐帧落盘?) then (是)
  :追加逐帧参数，上游每步写位姿与记忆状态;
  :解析 jsonl，写逐帧文件并登记 traj_path;
else (否)
  :降级为仅汇总指标;
endif
:汇总 SR / SPL / NE / nDTW 与资源峰值;
:回放接口读取任务号与行标识;
if (存在逐帧文件?) then (是)
  :还原俯视轨迹、第一视角、动作时间轴与记忆曲线;
else (否)
  :显示汇总指标与参考路径;
endif
stop
@enduml
\end{lstlisting}
\section{图~\ref{fig:flow05}：指标分析与诊断流程}

文件 \code{flow_05_analysis.png}。画法：从逐条结果到四条分析支路（对比、消融、显著性、长时曲线）与失败归因，最后汇入报告导出。

\begin{lstlisting}[style=navplain]
@startuml flow_05_analysis
skinparam defaultFontName "Noto Sans CJK SC"
title 指标分析与诊断流程
start
:读取逐条结果并聚合 SR / SPL / NE / nDTW;
fork
  :多任务对比矩阵与基线差值;
fork again
  :模态组合 × 记忆策略消融矩阵;
fork again
  :配对 bootstrap 与 95% 置信区间;
fork again
  :按步数分桶的长时曲线;
end fork
:九类失败归因自动分类与人工复核;
:生成 Markdown / CSV 报告并下载;
stop
@enduml
\end{lstlisting}
\section{图~\ref{fig:flow06}：交互式仿真子系统流程}

文件 \code{flow_06_sim.png}。画法：会话启动后 \code{fork} 出引擎循环、驾驶指令、录制三条并行支路，用 \code{end fork} 汇合到停止与资源释放。

\begin{lstlisting}[style=navplain]
@startuml flow_06_sim
skinparam defaultFontName "Noto Sans CJK SC"
title 交互式仿真子系统流程
start
:读取场景、起点、目标与参考路径;
:单会话检查并启动 Habitat 引擎线程;
fork
  :按节拍渲染第一视角与第三视角画中画;
  :MJPEG 分发并绘制俯视三线轨迹;
fork again
  :手动 / oracle / 模型动作写入队列并按节拍消费;
fork again
  :按 record_fps 采样 JPEG，停止后编码 mp4;
end fork
:输出步数、距离、NE 与 SPL 状态快照;
:空闲超时或用户停止后，停止引擎并释放资源;
stop
@enduml
\end{lstlisting}
\section{图~\ref{fig:flow07}：资产同步导入流程}

文件 \code{flow_07_backup_import.png}。画法：扫描→五类分类→签名比对→dry\_run 分支→写入→运行中结果增量同步。

\begin{lstlisting}[style=navplain]
@startuml flow_07_backup_import
skinparam defaultFontName "Noto Sans CJK SC"
title 资产同步导入流程
start
:读取绑定目录与自动同步开关;
:按五类扫描，排除大目录与超阈值路径;
:计算 source_key 与 signature;
if (与账本一致?) then (是)
  :跳过并保留已有记录;
else (否)
  if (dry_run 预览?) then (是)
    :只输出候选清单，不写库;
  else (否)
    :分类落库：数据集、资产索引、模型档案、历史任务;
  endif
endif
if (结果文件仍在写入?) then (是)
  :保持运行中并按条数增量刷新;
else (否)
  :转为已完成并补齐指标;
endif
:后台守护按间隔监视目录变化;
stop
@enduml
\end{lstlisting}
\section{界面截图：取图步骤}

\scriptsize
\begin{longtable}{@{}p{1.0cm}p{4.3cm}p{2.8cm}p{6.1cm}@{}}
\toprule
图号 & 文件名 & 路由 / 页面 & 取图步骤与画面要素 \\
\midrule
\endfirsthead
\toprule
图号 & 文件名 & 路由 / 页面 & 取图步骤与画面要素 \\
\midrule
\endhead
\bottomrule
\endfoot
\ref{fig:ui01} & \code{ui_01_overview.png} & \code{/overview} & 登录后首页；统计卡、队列、最近任务、趋势图与快捷入口 \\
\ref{fig:ui02} & \code{ui_02_dataset_list.png} & \code{/datasets} & 数据集列表；登记/扫描结果与场景、episode 统计 \\
\ref{fig:ui03} & \code{ui_03_observation_timeline.png} & \code{/datasets/:id} & 某 episode 的 RGB/深度/位姿/动作四轨时间轴与偏差标记 \\
\ref{fig:ui06} & \code{ui_06_task_wizard.png} & \code{/tasks} & 新建评测向导；数据集/切分/模型/模板选择与参数覆盖 \\
\ref{fig:ui07} & \code{ui_07_task_monitor.png} & \code{/tasks/:id} & 运行中任务；进度、当前 episode、日志与显存/延迟曲线 \\
\ref{fig:ui12} & \code{ui_12_analysis_compare.png} & \code{/analysis} & 选 2--3 个任务；指标矩阵与相对基线的差值列 \\
\ref{fig:ui14} & \code{ui_14_bootstrap.png} & \code{/analysis} & 两任务配对 bootstrap；均值差、95\% 置信区间与显著性 \\
\ref{fig:ui18} & \code{ui_18_sim_workbench.png} & \code{/sim} & 启动会话并走若干步；第一视角、俯视轨迹、控制与读数 \\
\ref{fig:ui19} & \code{ui_19_sim_policy.png} & \code{/sim} & 模型驾驶面板；加载/启用/暂停/停止与显存提示 \\
\ref{fig:ui23} & \code{ui_23_backup_import.png} & \code{/backup-import} & 绑定目录总览；五类候选数、扫描预览与自动同步开关 \\
\end{longtable}
\normalsize
。
% ============================================================

\chapter{插图绘制参考}

本说明书的插图统一放在 \code{figures/} 目录，文件名为小写英文与下划线；
结构图由 PlantUML 渲染为 PNG，界面截图在运行的系统中按路由截取。
图号、图题与交叉引用均由正文 \verb|\figph| 生成，增删图片后自动重排。

\section{通用出图约定}

\begin{itemize}[leftmargin=2em,itemsep=0.2em]
\item 结构图：复制下面的 \code{@startuml ... @enduml} 片段保存为同名 \code{.puml}，
  执行 \code{plantuml -tpng figures/<文件名去掉.png>.puml} 生成 PNG；中文渲染需本机中文字体。
\item 界面截图：启动 \code{python run.py --no-browser}，用 Chrome/Edge 在 1440×900 及以上分辨率截取 PNG。
\item 图片宽度建议不小于 1400 像素；正文按版心宽度 92\% 缩放，最大高度为版心高度 70\%。
\end{itemize}

\section{图~\ref{fig:arch01}：系统总体架构图}

文件 \code{arch_01_system_overview.png}，组件图。画法：按七层自上而下排列，前端三块（页面组、回放、仿真）→ 接口层三个组件 → 业务服务层五个组件 → 独立的仿真子系统包 → 适配器层 → 数据存储层 → 外部后端；用实线箭头标注 HTTP/JSON、SSE、MJPEG、子进程、文件读写等协议。

\begin{lstlisting}[style=navplain]
@startuml arch_01_system_overview
skinparam componentStyle rectangle
skinparam defaultFontName "Noto Sans CJK SC"
title 码极客 NavEval Studio V1.0 系统总体架构图
actor "用户 admin/operator/viewer" as User

package "浏览器前端 Vue 3" {
  [总览/数据/任务/分析页面组] as UI1
  [回放工作台 Canvas] as UI2
  [交互式仿真 MJPEG] as UI3
}
package "接口服务层 FastAPI" {
  [REST 路由 / Pydantic 校验] as API
  [Bearer + ?token 鉴权] as AUTH
  [SSE 事件流] as SSE
}
package "业务服务层" {
  [数据扫描/切分/质量校验] as SVC1
  [任务编排/调度器] as SVC2
  [指标聚合/归因/报告] as SVC3
  [回放解析/资产导入] as SVC4
  [后台守护 自动同步] as SVC5
}
package "仿真子系统（可选）" {
  [会话管理/动作队列] as SIM1
  [引擎线程 渲染循环] as SIM2
  [策略/录制/俯视绘制] as SIM3
}
package "适配器层" {
  [BaseAdapter 协议] as ADP
  [StreamVLNAdapter] as ADP1
  [MockAdapter] as ADP2
}
package "数据存储层" {
  [SQLite WAL eval.db] as DB
  [工作目录 storage/] as FS
}
package "外部推理与仿真后端" {
  [StreamVLN 推理子进程] as EXT1
  [habitat_sim 场景渲染] as EXT2
}
User --> UI1
UI1 --> API : HTTP / JSON
UI2 --> API : 回放 JSON + 帧图像
UI3 --> SSE : 事件流
UI3 --> API : MJPEG / ?token=
API --> AUTH
API --> SVC1
API --> SVC2
API --> SVC3
API --> SVC4
SVC2 --> SSE : 进度/采样/日志
SVC2 --> ADP
ADP <|-- ADP1
ADP <|-- ADP2
ADP1 --> EXT1 : 子进程 + --dump_frames
SVC4 --> SIM1
SIM1 --> SIM2
SIM1 --> SIM3
SIM2 --> EXT2
SIM3 --> EXT2
SVC1 --> DB
SVC2 --> DB
SVC3 --> DB
SVC4 --> DB
SVC5 --> DB
SVC5 --> FS : 备份库目录监视
SVC2 --> FS : 快照/日志/逐帧 jsonl
SIM3 --> FS : 录制帧与 mp4
@enduml
\end{lstlisting}
\section{图~\ref{fig:flow01}：评测任务执行总流程}

文件 \code{flow_01_eval_task.png}，活动图，软著材料的核心流程图。画法：用泳道区分前端、接口服务层、调度器、适配器/子进程、数据存储，覆盖提交、校验、排队、执行、事件回写、聚合六个阶段。

\begin{lstlisting}[style=navplain]
@startuml flow_01_eval_task
skinparam defaultFontName "Noto Sans CJK SC"
title 评测任务执行总流程
|前端|
start
:选择数据集/切分/模型/模板并覆盖参数;
:提交创建任务;
|接口服务层|
:参数校验 + 因果性白名单预检;
if (校验通过?) then (否)
  :返回 1001 / 1006 与具体字段;
  |前端|
  :就地提示并保留表单;
  stop
else (是)
endif
|数据存储|
:写入 eval_task(queued) 与配置快照;
|调度器|
:入队，等待并发额度与 GPU 互斥;
:取出任务，创建工作目录;
:写 config.json 与启动命令;
|适配器/子进程|
:prepare 校验路径与后端能力;
:拉起推理子进程（可选 --dump_frames）;
repeat
  :逐条执行 episode;
  :产出进度/资源/日志/逐帧记录;
  |数据存储|
  :按行解析写 task_episode 与 frames/*.jsonl;
  |前端|
  :SSE 更新进度条/曲线/日志;
repeat while (还有 episode?) is (是)
|接口服务层|
:聚合指标、生成回放索引与报告索引;
:任务置为 succeeded，推送 done 事件;
|前端|
:展示结果，支持对比/回放/导出;
stop
@enduml
\end{lstlisting}
\section{图~\ref{fig:flow02}：数据与配置管理子系统流程}

文件 \code{flow_02_data_manage.png}。画法：一条主干（登记→扫描→索引→校验→切分→导出），在校验处分支展示问题清单，在导出处强调时间前缀二次校验。

\begin{lstlisting}[style=navplain]
@startuml flow_02_data_manage
skinparam defaultFontName "Noto Sans CJK SC"
title 数据与配置管理子系统流程
start
:登记数据集并校验根路径白名单;
if (路径合法且存在?) then (否)
  :状态置为 missing 并提示修正;
  stop
else (是)
endif
:扫描 episode 并写入索引，状态置为 ready;
:多源观测对齐与质量校验;
if (存在问题?) then (是)
  :输出问题清单与涉及 episode;
else (否)
  :标记校验通过;
endif
:生成 full / smoke / subset 切分并保存条件;
:因果切片导出与时间前缀二次校验;
:写出 jsonl 并返回下载链接;
stop
@enduml
\end{lstlisting}
\section{图~\ref{fig:flow03}：任务调度子系统流程}

文件 \code{flow_03_scheduler.png}。画法：循环从队列取任务，突出并发额度、GPU 互斥、episode 边界上的暂停/取消检查与失败重跑。

\begin{lstlisting}[style=navplain]
@startuml flow_03_scheduler
skinparam defaultFontName "Noto Sans CJK SC"
title 任务调度子系统流程
start
while (队列非空?) is (是)
  :按优先级取队首任务;
  if (并发额度未满且 GPU 可用?) then (否)
    :留在队列等待下一轮;
  else (是)
    :标记 running 并占用 GPU;
    :prepare 调用适配器;
    while (还有 episode?) is (是)
      :执行一条 episode，写结果并推送事件;
      if (收到暂停?) then (是)
        :在 episode 边界置 paused 并等待恢复;
      else (否)
      endif
      if (收到取消?) then (是)
        :终止子进程，置 cancelled 并释放资源;
        stop
      else (否)
      endif
    endwhile
    :聚合指标并释放并发额度与 GPU;
  endif
endwhile (否)
:队列空闲，等待新任务或重跑;
stop
@enduml
\end{lstlisting}
\section{图~\ref{fig:flow04}：推理适配器与逐帧回放流程}

文件 \code{flow_04_adapter_replay.png}。画法：从 prepare 探测上游能力开始，分“支持/不支持逐帧”两支，最终汇合到回放接口与四区工作台。

\begin{lstlisting}[style=navplain]
@startuml flow_04_adapter_replay
skinparam defaultFontName "Noto Sans CJK SC"
title 推理适配器与逐帧回放流程
start
:读取配置快照，prepare 校验路径与权重;
if (上游支持逐帧落盘?) then (是)
  :追加逐帧参数，上游每步写位姿与记忆状态;
  :解析 jsonl，写逐帧文件并登记 traj_path;
else (否)
  :降级为仅汇总指标;
endif
:汇总 SR / SPL / NE / nDTW 与资源峰值;
:回放接口读取任务号与行标识;
if (存在逐帧文件?) then (是)
  :还原俯视轨迹、第一视角、动作时间轴与记忆曲线;
else (否)
  :显示汇总指标与参考路径;
endif
stop
@enduml
\end{lstlisting}
\section{图~\ref{fig:flow05}：指标分析与诊断流程}

文件 \code{flow_05_analysis.png}。画法：从逐条结果到四条分析支路（对比、消融、显著性、长时曲线）与失败归因，最后汇入报告导出。

\begin{lstlisting}[style=navplain]
@startuml flow_05_analysis
skinparam defaultFontName "Noto Sans CJK SC"
title 指标分析与诊断流程
start
:读取逐条结果并聚合 SR / SPL / NE / nDTW;
fork
  :多任务对比矩阵与基线差值;
fork again
  :模态组合 × 记忆策略消融矩阵;
fork again
  :配对 bootstrap 与 95% 置信区间;
fork again
  :按步数分桶的长时曲线;
end fork
:九类失败归因自动分类与人工复核;
:生成 Markdown / CSV 报告并下载;
stop
@enduml
\end{lstlisting}
\section{图~\ref{fig:flow06}：交互式仿真子系统流程}

文件 \code{flow_06_sim.png}。画法：会话启动后 \code{fork} 出引擎循环、驾驶指令、录制三条并行支路，用 \code{end fork} 汇合到停止与资源释放。

\begin{lstlisting}[style=navplain]
@startuml flow_06_sim
skinparam defaultFontName "Noto Sans CJK SC"
title 交互式仿真子系统流程
start
:读取场景、起点、目标与参考路径;
:单会话检查并启动 Habitat 引擎线程;
fork
  :按节拍渲染第一视角与第三视角画中画;
  :MJPEG 分发并绘制俯视三线轨迹;
fork again
  :手动 / oracle / 模型动作写入队列并按节拍消费;
fork again
  :按 record_fps 采样 JPEG，停止后编码 mp4;
end fork
:输出步数、距离、NE 与 SPL 状态快照;
:空闲超时或用户停止后，停止引擎并释放资源;
stop
@enduml
\end{lstlisting}
\section{图~\ref{fig:flow07}：资产同步导入流程}

文件 \code{flow_07_backup_import.png}。画法：扫描→五类分类→签名比对→dry\_run 分支→写入→运行中结果增量同步。

\begin{lstlisting}[style=navplain]
@startuml flow_07_backup_import
skinparam defaultFontName "Noto Sans CJK SC"
title 资产同步导入流程
start
:读取绑定目录与自动同步开关;
:按五类扫描，排除大目录与超阈值路径;
:计算 source_key 与 signature;
if (与账本一致?) then (是)
  :跳过并保留已有记录;
else (否)
  if (dry_run 预览?) then (是)
    :只输出候选清单，不写库;
  else (否)
    :分类落库：数据集、资产索引、模型档案、历史任务;
  endif
endif
if (结果文件仍在写入?) then (是)
  :保持运行中并按条数增量刷新;
else (否)
  :转为已完成并补齐指标;
endif
:后台守护按间隔监视目录变化;
stop
@enduml
\end{lstlisting}
\section{界面截图：取图步骤}

\scriptsize
\begin{longtable}{@{}p{1.0cm}p{4.3cm}p{2.8cm}p{6.1cm}@{}}
\toprule
图号 & 文件名 & 路由 / 页面 & 取图步骤与画面要素 \\
\midrule
\endfirsthead
\toprule
图号 & 文件名 & 路由 / 页面 & 取图步骤与画面要素 \\
\midrule
\endhead
\bottomrule
\endfoot
\ref{fig:ui01} & \code{ui_01_overview.png} & \code{/overview} & 登录后首页；统计卡、队列、最近任务、趋势图与快捷入口 \\
\ref{fig:ui02} & \code{ui_02_dataset_list.png} & \code{/datasets} & 数据集列表；登记/扫描结果与场景、episode 统计 \\
\ref{fig:ui03} & \code{ui_03_observation_timeline.png} & \code{/datasets/:id} & 某 episode 的 RGB/深度/位姿/动作四轨时间轴与偏差标记 \\
\ref{fig:ui06} & \code{ui_06_task_wizard.png} & \code{/tasks} & 新建评测向导；数据集/切分/模型/模板选择与参数覆盖 \\
\ref{fig:ui07} & \code{ui_07_task_monitor.png} & \code{/tasks/:id} & 运行中任务；进度、当前 episode、日志与显存/延迟曲线 \\
\ref{fig:ui12} & \code{ui_12_analysis_compare.png} & \code{/analysis} & 选 2--3 个任务；指标矩阵与相对基线的差值列 \\
\ref{fig:ui14} & \code{ui_14_bootstrap.png} & \code{/analysis} & 两任务配对 bootstrap；均值差、95\% 置信区间与显著性 \\
\ref{fig:ui18} & \code{ui_18_sim_workbench.png} & \code{/sim} & 启动会话并走若干步；第一视角、俯视轨迹、控制与读数 \\
\ref{fig:ui19} & \code{ui_19_sim_policy.png} & \code{/sim} & 模型驾驶面板；加载/启用/暂停/停止与显存提示 \\
\ref{fig:ui23} & \code{ui_23_backup_import.png} & \code{/backup-import} & 绑定目录总览；五类候选数、扫描预览与自动同步开关 \\
\end{longtable}
\normalsize
。
% ============================================================

\chapter{插图绘制参考}

本说明书的插图统一放在 \code{figures/} 目录，文件名为小写英文与下划线；
结构图由 PlantUML 渲染为 PNG，界面截图在运行的系统中按路由截取。
图号、图题与交叉引用均由正文 \verb|\figph| 生成，增删图片后自动重排。

\section{通用出图约定}

\begin{itemize}[leftmargin=2em,itemsep=0.2em]
\item 结构图：复制下面的 \code{@startuml ... @enduml} 片段保存为同名 \code{.puml}，
  执行 \code{plantuml -tpng figures/<文件名去掉.png>.puml} 生成 PNG；中文渲染需本机中文字体。
\item 界面截图：启动 \code{python run.py --no-browser}，用 Chrome/Edge 在 1440×900 及以上分辨率截取 PNG。
\item 图片宽度建议不小于 1400 像素；正文按版心宽度 92\% 缩放，最大高度为版心高度 70\%。
\end{itemize}

\section{图~\ref{fig:arch01}：系统总体架构图}

文件 \code{arch_01_system_overview.png}，组件图。画法：按七层自上而下排列，前端三块（页面组、回放、仿真）→ 接口层三个组件 → 业务服务层五个组件 → 独立的仿真子系统包 → 适配器层 → 数据存储层 → 外部后端；用实线箭头标注 HTTP/JSON、SSE、MJPEG、子进程、文件读写等协议。

\begin{lstlisting}[style=navplain]
@startuml arch_01_system_overview
skinparam componentStyle rectangle
skinparam defaultFontName "Noto Sans CJK SC"
title 码极客 NavEval Studio V1.0 系统总体架构图
actor "用户 admin/operator/viewer" as User

package "浏览器前端 Vue 3" {
  [总览/数据/任务/分析页面组] as UI1
  [回放工作台 Canvas] as UI2
  [交互式仿真 MJPEG] as UI3
}
package "接口服务层 FastAPI" {
  [REST 路由 / Pydantic 校验] as API
  [Bearer + ?token 鉴权] as AUTH
  [SSE 事件流] as SSE
}
package "业务服务层" {
  [数据扫描/切分/质量校验] as SVC1
  [任务编排/调度器] as SVC2
  [指标聚合/归因/报告] as SVC3
  [回放解析/资产导入] as SVC4
  [后台守护 自动同步] as SVC5
}
package "仿真子系统（可选）" {
  [会话管理/动作队列] as SIM1
  [引擎线程 渲染循环] as SIM2
  [策略/录制/俯视绘制] as SIM3
}
package "适配器层" {
  [BaseAdapter 协议] as ADP
  [StreamVLNAdapter] as ADP1
  [MockAdapter] as ADP2
}
package "数据存储层" {
  [SQLite WAL eval.db] as DB
  [工作目录 storage/] as FS
}
package "外部推理与仿真后端" {
  [StreamVLN 推理子进程] as EXT1
  [habitat_sim 场景渲染] as EXT2
}
User --> UI1
UI1 --> API : HTTP / JSON
UI2 --> API : 回放 JSON + 帧图像
UI3 --> SSE : 事件流
UI3 --> API : MJPEG / ?token=
API --> AUTH
API --> SVC1
API --> SVC2
API --> SVC3
API --> SVC4
SVC2 --> SSE : 进度/采样/日志
SVC2 --> ADP
ADP <|-- ADP1
ADP <|-- ADP2
ADP1 --> EXT1 : 子进程 + --dump_frames
SVC4 --> SIM1
SIM1 --> SIM2
SIM1 --> SIM3
SIM2 --> EXT2
SIM3 --> EXT2
SVC1 --> DB
SVC2 --> DB
SVC3 --> DB
SVC4 --> DB
SVC5 --> DB
SVC5 --> FS : 备份库目录监视
SVC2 --> FS : 快照/日志/逐帧 jsonl
SIM3 --> FS : 录制帧与 mp4
@enduml
\end{lstlisting}
\section{图~\ref{fig:flow01}：评测任务执行总流程}

文件 \code{flow_01_eval_task.png}，活动图，软著材料的核心流程图。画法：用泳道区分前端、接口服务层、调度器、适配器/子进程、数据存储，覆盖提交、校验、排队、执行、事件回写、聚合六个阶段。

\begin{lstlisting}[style=navplain]
@startuml flow_01_eval_task
skinparam defaultFontName "Noto Sans CJK SC"
title 评测任务执行总流程
|前端|
start
:选择数据集/切分/模型/模板并覆盖参数;
:提交创建任务;
|接口服务层|
:参数校验 + 因果性白名单预检;
if (校验通过?) then (否)
  :返回 1001 / 1006 与具体字段;
  |前端|
  :就地提示并保留表单;
  stop
else (是)
endif
|数据存储|
:写入 eval_task(queued) 与配置快照;
|调度器|
:入队，等待并发额度与 GPU 互斥;
:取出任务，创建工作目录;
:写 config.json 与启动命令;
|适配器/子进程|
:prepare 校验路径与后端能力;
:拉起推理子进程（可选 --dump_frames）;
repeat
  :逐条执行 episode;
  :产出进度/资源/日志/逐帧记录;
  |数据存储|
  :按行解析写 task_episode 与 frames/*.jsonl;
  |前端|
  :SSE 更新进度条/曲线/日志;
repeat while (还有 episode?) is (是)
|接口服务层|
:聚合指标、生成回放索引与报告索引;
:任务置为 succeeded，推送 done 事件;
|前端|
:展示结果，支持对比/回放/导出;
stop
@enduml
\end{lstlisting}
\section{图~\ref{fig:flow02}：数据与配置管理子系统流程}

文件 \code{flow_02_data_manage.png}。画法：一条主干（登记→扫描→索引→校验→切分→导出），在校验处分支展示问题清单，在导出处强调时间前缀二次校验。

\begin{lstlisting}[style=navplain]
@startuml flow_02_data_manage
skinparam defaultFontName "Noto Sans CJK SC"
title 数据与配置管理子系统流程
start
:登记数据集并校验根路径白名单;
if (路径合法且存在?) then (否)
  :状态置为 missing 并提示修正;
  stop
else (是)
endif
:扫描 episode 并写入索引，状态置为 ready;
:多源观测对齐与质量校验;
if (存在问题?) then (是)
  :输出问题清单与涉及 episode;
else (否)
  :标记校验通过;
endif
:生成 full / smoke / subset 切分并保存条件;
:因果切片导出与时间前缀二次校验;
:写出 jsonl 并返回下载链接;
stop
@enduml
\end{lstlisting}
\section{图~\ref{fig:flow03}：任务调度子系统流程}

文件 \code{flow_03_scheduler.png}。画法：循环从队列取任务，突出并发额度、GPU 互斥、episode 边界上的暂停/取消检查与失败重跑。

\begin{lstlisting}[style=navplain]
@startuml flow_03_scheduler
skinparam defaultFontName "Noto Sans CJK SC"
title 任务调度子系统流程
start
while (队列非空?) is (是)
  :按优先级取队首任务;
  if (并发额度未满且 GPU 可用?) then (否)
    :留在队列等待下一轮;
  else (是)
    :标记 running 并占用 GPU;
    :prepare 调用适配器;
    while (还有 episode?) is (是)
      :执行一条 episode，写结果并推送事件;
      if (收到暂停?) then (是)
        :在 episode 边界置 paused 并等待恢复;
      else (否)
      endif
      if (收到取消?) then (是)
        :终止子进程，置 cancelled 并释放资源;
        stop
      else (否)
      endif
    endwhile
    :聚合指标并释放并发额度与 GPU;
  endif
endwhile (否)
:队列空闲，等待新任务或重跑;
stop
@enduml
\end{lstlisting}
\section{图~\ref{fig:flow04}：推理适配器与逐帧回放流程}

文件 \code{flow_04_adapter_replay.png}。画法：从 prepare 探测上游能力开始，分“支持/不支持逐帧”两支，最终汇合到回放接口与四区工作台。

\begin{lstlisting}[style=navplain]
@startuml flow_04_adapter_replay
skinparam defaultFontName "Noto Sans CJK SC"
title 推理适配器与逐帧回放流程
start
:读取配置快照，prepare 校验路径与权重;
if (上游支持逐帧落盘?) then (是)
  :追加逐帧参数，上游每步写位姿与记忆状态;
  :解析 jsonl，写逐帧文件并登记 traj_path;
else (否)
  :降级为仅汇总指标;
endif
:汇总 SR / SPL / NE / nDTW 与资源峰值;
:回放接口读取任务号与行标识;
if (存在逐帧文件?) then (是)
  :还原俯视轨迹、第一视角、动作时间轴与记忆曲线;
else (否)
  :显示汇总指标与参考路径;
endif
stop
@enduml
\end{lstlisting}
\section{图~\ref{fig:flow05}：指标分析与诊断流程}

文件 \code{flow_05_analysis.png}。画法：从逐条结果到四条分析支路（对比、消融、显著性、长时曲线）与失败归因，最后汇入报告导出。

\begin{lstlisting}[style=navplain]
@startuml flow_05_analysis
skinparam defaultFontName "Noto Sans CJK SC"
title 指标分析与诊断流程
start
:读取逐条结果并聚合 SR / SPL / NE / nDTW;
fork
  :多任务对比矩阵与基线差值;
fork again
  :模态组合 × 记忆策略消融矩阵;
fork again
  :配对 bootstrap 与 95% 置信区间;
fork again
  :按步数分桶的长时曲线;
end fork
:九类失败归因自动分类与人工复核;
:生成 Markdown / CSV 报告并下载;
stop
@enduml
\end{lstlisting}
\section{图~\ref{fig:flow06}：交互式仿真子系统流程}

文件 \code{flow_06_sim.png}。画法：会话启动后 \code{fork} 出引擎循环、驾驶指令、录制三条并行支路，用 \code{end fork} 汇合到停止与资源释放。

\begin{lstlisting}[style=navplain]
@startuml flow_06_sim
skinparam defaultFontName "Noto Sans CJK SC"
title 交互式仿真子系统流程
start
:读取场景、起点、目标与参考路径;
:单会话检查并启动 Habitat 引擎线程;
fork
  :按节拍渲染第一视角与第三视角画中画;
  :MJPEG 分发并绘制俯视三线轨迹;
fork again
  :手动 / oracle / 模型动作写入队列并按节拍消费;
fork again
  :按 record_fps 采样 JPEG，停止后编码 mp4;
end fork
:输出步数、距离、NE 与 SPL 状态快照;
:空闲超时或用户停止后，停止引擎并释放资源;
stop
@enduml
\end{lstlisting}
\section{图~\ref{fig:flow07}：资产同步导入流程}

文件 \code{flow_07_backup_import.png}。画法：扫描→五类分类→签名比对→dry\_run 分支→写入→运行中结果增量同步。

\begin{lstlisting}[style=navplain]
@startuml flow_07_backup_import
skinparam defaultFontName "Noto Sans CJK SC"
title 资产同步导入流程
start
:读取绑定目录与自动同步开关;
:按五类扫描，排除大目录与超阈值路径;
:计算 source_key 与 signature;
if (与账本一致?) then (是)
  :跳过并保留已有记录;
else (否)
  if (dry_run 预览?) then (是)
    :只输出候选清单，不写库;
  else (否)
    :分类落库：数据集、资产索引、模型档案、历史任务;
  endif
endif
if (结果文件仍在写入?) then (是)
  :保持运行中并按条数增量刷新;
else (否)
  :转为已完成并补齐指标;
endif
:后台守护按间隔监视目录变化;
stop
@enduml
\end{lstlisting}
\section{界面截图：取图步骤}

\scriptsize
\begin{longtable}{@{}p{1.0cm}p{4.3cm}p{2.8cm}p{6.1cm}@{}}
\toprule
图号 & 文件名 & 路由 / 页面 & 取图步骤与画面要素 \\
\midrule
\endfirsthead
\toprule
图号 & 文件名 & 路由 / 页面 & 取图步骤与画面要素 \\
\midrule
\endhead
\bottomrule
\endfoot
\ref{fig:ui01} & \code{ui_01_overview.png} & \code{/overview} & 登录后首页；统计卡、队列、最近任务、趋势图与快捷入口 \\
\ref{fig:ui02} & \code{ui_02_dataset_list.png} & \code{/datasets} & 数据集列表；登记/扫描结果与场景、episode 统计 \\
\ref{fig:ui03} & \code{ui_03_observation_timeline.png} & \code{/datasets/:id} & 某 episode 的 RGB/深度/位姿/动作四轨时间轴与偏差标记 \\
\ref{fig:ui06} & \code{ui_06_task_wizard.png} & \code{/tasks} & 新建评测向导；数据集/切分/模型/模板选择与参数覆盖 \\
\ref{fig:ui07} & \code{ui_07_task_monitor.png} & \code{/tasks/:id} & 运行中任务；进度、当前 episode、日志与显存/延迟曲线 \\
\ref{fig:ui12} & \code{ui_12_analysis_compare.png} & \code{/analysis} & 选 2--3 个任务；指标矩阵与相对基线的差值列 \\
\ref{fig:ui14} & \code{ui_14_bootstrap.png} & \code{/analysis} & 两任务配对 bootstrap；均值差、95\% 置信区间与显著性 \\
\ref{fig:ui18} & \code{ui_18_sim_workbench.png} & \code{/sim} & 启动会话并走若干步；第一视角、俯视轨迹、控制与读数 \\
\ref{fig:ui19} & \code{ui_19_sim_policy.png} & \code{/sim} & 模型驾驶面板；加载/启用/暂停/停止与显存提示 \\
\ref{fig:ui23} & \code{ui_23_backup_import.png} & \code{/backup-import} & 绑定目录总览；五类候选数、扫描预览与自动同步开关 \\
\end{longtable}
\normalsize
